% ============================================================================
% CCSE-CS for EMNLP 2026 - Main Submission Version
% Limited to 8 pages (excluding references)
% Appendix and extended examples moved to supplementary material
% ============================================================================

\documentclass{article}

% EMNLP 2026 packages
\usepackage[utf8]{inputenc}
\usepackage[T1]{fontenc}
\usepackage{hyperref}
\usepackage{url}
\usepackage{booktabs}
\usepackage{amsfonts}
\usepackage{nicefrac}
\usepackage{microtype}
\usepackage{cite}
\usepackage{amsmath,amssymb,amsthm}
\usepackage{tikz}
\usepackage{pgfplots}
\pgfplotsset{compat=1.17}

% Chinese support
\usepackage[UTF8]{ctex}

% Page setup (EMNLP format)
\usepackage{geometry}
\geometry{letterpaper,tmargin=1in,bmargin=1in,lmargin=1in,rmargin=1in}

% Remove all anonymity-identifying information
\title{CCSE-CS: A Culturally-Aware Empathetic Dialogue Dataset for Chinese Customer Service}

\author{
  Anonymous Authors \\
  Anonymous Institution
}

\begin{document}

\maketitle

\begin{abstract}
Customer service dialogue systems require deep cultural understanding and empathetic communication skills, especially in high-context cultures like China where face-saving (面子) and indirect communication patterns play crucial roles. While existing empathetic dialogue datasets (ESConv, CSC) have advanced English-centric customer service AI, they lack (1) Chinese cultural grounding, (2) fine-grained strategy taxonomies, and (3) rigorous multi-dimensional evaluation protocols.

We present \textbf{CCSE-CS} (\textit{Customer Service Culturally-Aware Strategies for Empathy}), a high-quality Chinese customer service dialogue dataset with 522 dialogues and 3,285 turns across four domains (e-commerce, banking, telecommunications, healthcare). Our key contributions include:

\begin{itemize}
  \item A \textbf{four-dimensional cultural characterization framework} (relationship orientation, face sensitivity, euphemism level, conflict intensity) enabling systematic study of cultural factors in empathetic dialogue, validated through ablation studies (+0.70 performance gain).

  \item A \textbf{6-category $\times$ 18-substrategy taxonomy} derived from cross-cultural psychology and customer service standards, with explicit evidence chains linking each strategy to existing literature or cultural adaptations. Ablation studies validate fine-grained strategies over coarse-grained taxonomies (+0.37 performance gain).

  \item A \textbf{four-dimensional evaluation protocol} (Empathy-Appropriateness, Policy-Compliance, Cultural-Fit, Naturalness) extending recent LLM-as-a-judge work with customer-service-specific dimensions, achieving Pearson \(r=0.969\) with human experts.

  \item A \textbf{cost-effective annotation pipeline} combining GPT-4 and human experts, achieving 83.3\% cost reduction (\$400 to \$75) while maintaining quality (Fleiss' Kappa=0.82).
\end{itemize}

Experiments show that models trained on CCSE-CS achieve 6.9\%+ improvement in overall ratings when scaling from 100 to 522 dialogues, with monotonic performance gains across all dimensions. Ablation studies validate the necessity of fine-grained strategies over coarse-grained taxonomies (e.g., splitting Emotional Management into 5 sub-strategies).

\textbf{中文摘要：} 客服对话系统需要深刻的文化理解和共情沟通能力，尤其在中国等高语境文化中，"面子"维护和委婉表达至关重要。现有共情对话数据集（ESConv、CSC）虽推动了英文客服AI发展，但缺乏(1)中国文化扎根、(2)细粒度策略分类、(3)文化感知评测协议。本文提出 \textbf{CCSE-CS}（中文客服文化感知共情对话数据集），包含522段对话和3,285轮对话，覆盖电商/银行/电信/医疗四领域。核心贡献：(1) \textbf{四维文化表征框架}（关系取向/面子敏感度/委婉度/冲突强度），消融验证文化因素贡献+0.70；(2) \textbf{6类18子策略本体}（继承4个/升级6个/创新8个），消融验证细粒度优于粗粒度（+0.37）；(3) \textbf{四维评测协议}（共情适时性/政策合规/文化适切/自然度），r=0.969；(4) \textbf{成本有效标注管线}，节省83.3\%成本（Kappa=0.82）。数据量从100条扩展到522条时，整体评分提升6.9\%+，所有维度性能单调递增。
\end{abstract}

% [Content from main.tex will be included here, focusing on:
% - Introduction (1 page)
% - Related Work (1 page, compressed)
% - Method (2 pages)
% - Experiments (3 pages, focus on key results)
% - Discussion & Limitations (1 page)
%
% Note: Full appendix, strategy details, and extended examples
% are moved to supplementary_material.pdf]

% [COPY RELEVANT SECTIONS FROM main.tex HERE]

% ============================================================================
% References
% ============================================================================

\bibliographystyle{acmtrans}
\bibliography{references}

\end{document}

% ============================================================================
% Submission Notes:
% - This version fits within 8-page EMNLP limit
% - Supplementary material (appendix.pdf) contains:
%   * Full strategy derivation evidence chains
%   * Quality filter examples
%   * Cultural persona profiles
%   * Extended error analysis
%   * Per-strategy IAA details
% ============================================================================
